\documentclass[usenames,dvipsnames, 18pt, compress, aspectratio=169]{beamer}

% can be compiled by xelatex -shell-escape presentation.tex
% lualatex -shell-escape presentation.tex

\usetheme[]{metropolis}

\usepackage[utf8]{inputenc}
\usepackage[russian, english]{babel}
\usepackage{booktabs}
\usepackage[scale=2]{ccicons}
\usepackage{listings}
\usepackage{marvosym}
\usepackage{color}
\usepackage{xcolor}
\usepackage[document]{ragged2e}
\usepackage[export]{adjustbox}
\usepackage{fontawesome5}
\usepackage{enumitem}
\usepackage{minted}
\usemintedstyle{tango}
\usepackage[normalem]{ulem}
\usepackage{tikz}
\usetikzlibrary{patterns}
\usetikzlibrary{mindmap}
\usetikzlibrary{shapes.misc, fit}
\usetikzlibrary{spy, decorations, decorations.pathmorphing, decorations.pathreplacing, backgrounds, decorations.text}
\usepackage{graphicx}
\usepackage{eso-pic}
\usepackage{verbatim}
\usepackage{smartdiagram}
\usesmartdiagramlibrary{additions}
\usetikzlibrary{trees}
\usepackage{datetime}
\usepackage{hyperref}
\usepackage{forloop}
\usepackage{copyrightbox}
\usepackage{csquotes}
\usepackage{hyperref}
\usepackage{pgfplots}
\usetikzlibrary{tikzmark}
\usetikzlibrary{arrows.meta}

\usepackage{tcolorbox}
\usepackage{tabularx}
\usepackage{array}
\usepackage{colortbl}
\tcbuselibrary{skins}

\usetikzlibrary{shapes,arrows,positioning}
\graphicspath{{images/}}
%\newfontfamily{\FA}{FontAwesome}

\def\social{{\faMastodon}}
\def\github{{\faGithub}}
\def\email{{\faEnvelope}}

\renewcommand{\ttdefault}{pcr}
\setmonofont{FiraCode-VF}
%\newfontfamily{\ttfamily}{FiraCode-VF}

\usefonttheme{professionalfonts} % using non standard fonts for beamer
\usefonttheme{serif} % default family is serif
\usepackage{fontspec}
\setmainfont{Liberation Sans}
\newfontfamily\ExtraLight{Liberation Sans}
\newfontfamily\Light{Liberation Sans}
\newfontfamily\Book{Liberation Sans}
\newfontfamily\Medium{Liberation Sans}

\makeatletter
\newcommand\HUGE{\@setfontsize\Huge{32}{41}}
\makeatother

\newcommand\AtPagemyUpperLeft[1]{\AtPageLowerLeft{%
\put(\LenToUnit{0.85\paperwidth},\LenToUnit{0.05\paperheight}){#1}}}

\newcommand\AtPagemyUpperTop[1]{\AtPageLowerLeft{%
\put(\LenToUnit{0.42\paperwidth},\LenToUnit{0.90\paperheight}){#1}}}

\renewcommand{\ULthickness}{2.0pt}

\definecolor{links}{HTML}{0099FF}
\hypersetup{colorlinks, linkcolor=, urlcolor=links}
\definecolor{greenGood}{HTML}{99FF99}
\definecolor{redBad}{HTML}{FF9980}
\definecolor{title}{HTML}{ee0000}

\setbeamerfont{section title}{family=\Book, size=\Huge, shape=\normalfont}
\setbeamerfont{frametitle}{family=\Book, size=\large, shape=\normalfont}
\setbeamerfont{title}{family=\Book, size=\Large, shape=\normalfont}
\setbeamerfont{subtitle}{size=\small}
\setbeamerfont{author}{family=\ExtraLight, size=\footnotesize}

\definecolor{cec1d24}{RGB}{236,29,36}
\definecolor{cffffff}{RGB}{255,255,255}

\setbeamertemplate{navigation symbols}{}
\beamertemplatenavigationsymbolsempty
\pagenumbering{gobble}

\setbeamertemplate{title page}
{

  \vspace*{2.1cm}
  \hspace{6.0cm}
  \begin{minipage}[b][\paperheight]{0.6\textwidth}
  \begin{center}

    \ifx\inserttitle\@empty\else
    {{% \inserttitle is nonempty
      \raggedright%
      %\linespread{1.0}%
      \usebeamerfont{title}%
      \usebeamercolor[fg]{title}%
      %\vspace*{1.3em}
      \if@noSmallCapitals%
        \inserttitle%
      \else%
        \scshape{\color{title} \textbf{\begin{flushleft}\inserttitle\end{flushleft}}}%
      \fi%
      \vspace*{0.3em}
    }}
    \fi

    \vspace*{0.5em}%

    \ifx\insertsubtitle\@empty\else
    {{% \insertsubtitle is nonempty
      \usebeamerfont{subtitle}%
      \usebeamercolor[fg]{subtitle}%
      {\color{black} \insertsubtitle}%
      \vspace*{3.0em}%
    }}
    \fi

    \vspace*{1.0em}%

    \usebeamerfont{author}%
    \usebeamercolor[fg]{author}%
    {\begin{flushleft}\color{black} \insertauthor\end{flushleft}}%

    %\vspace*{1.5em}
    \fontsize{8pt}{10}\selectfont
    {\begin{flushleft}\color{black} 15-12-2022\end{flushleft}}%

    \vfill
    \vspace*{2em}
  \end{center}
  \end{minipage}
}

\setbeamertemplate{section page}
{
  \vspace{2em}
  \centering
  \begin{minipage}{22em}
    \usebeamercolor[fg]{section title}
    \usebeamerfont{section title}
    {\color{black} \insertsectionhead\\[-1ex]}
  \end{minipage}
  \par
}

\setbeamertemplate{footline}
{
\begin{beamercolorbox}[wd=\textwidth,ht=3ex,dp=3ex,leftskip=0.3cm,rightskip=0.3cm]{structure}
  \usebeamerfont{page number in head/foot}
  \insertframenumber
\end{beamercolorbox}
}


\title{Demanding the impossible: \\ rigorous database benchmarking}
\subtitle{}
\date{\today}
\author{DMITRII DOLGOV}
\institute{}

\begin{document}
{
  \usebackgroundtemplate{\includegraphics[width=\paperwidth]{title.png}}%
  \fontsize{17pt}{18}\selectfont
  \maketitle
}

\AddToShipoutPictureBG{
  \AtPagemyUpperLeft{{\includegraphics[width=2.0cm,keepaspectratio]{logo.png}}}
}

%\setbeamertemplate{background canvas}{}

\fontsize{17pt}{18}\selectfont

% party tricks references

%{
%\setbeamercolor{background canvas}{bg=black}
%\begin{frame}
    %\frametitle{}
    %\begin{center}
        %\includegraphics[width=1.0\textwidth,center]{nethack.png}
    %\end{center}
%\end{frame}
%\note{Originally about resources, but now not only}
%}

\setbeamertemplate{background canvas}{}

\section{Benchmarking model}

\begin{frame}
    \frametitle{}
    \begin{center}
        Motivation: Choose your fighter

        https://github.com/cmu-db/benchbase
        https://github.com/akopytov/sysbench
        https://www.postgresql.org/docs/current/pgbench.html
        https://github.com/brianfrankcooper/YCSB
        https://github.com/TPC-Council/HammerDB

        Replicated prod workload
    \end{center}
\end{frame}

\begin{frame}[fragile]{}
    \frametitle{}
    \begin{center}
        Motivation

		\begin{overprint}[14.5cm]
        \onslide<1>
        \begin{minted}[fontsize=\large]{shell-session}
latency average = 0.011 ms
latency stddev = 0.002 ms
tps = 89357.630697 (without initial connection time)
        \end{minted}

        \onslide<2>
        \begin{minted}[fontsize=\large]{shell-session}
latency average = 0.011 ms
latency stddev = 0.002 ms
tps = 89357.630697 (without initial connection time)

latency average = 0.014 ms
latency stddev = 0.023 ms
tps = 67107.536620 (without initial connection time)
        \end{minted}
		\end{overprint}

    \end{center}
\end{frame}

\begin{frame}
    \frametitle{}
    \begin{center}
        \copyrightbox[b]
            {\includegraphics[width=0.8\textwidth,center]{lorenz.png}}
            {
                The phase space plot of the Lorenz attractor,
                \\ {\scriptsize
                    \href{https://www.researchgate.net/publication/339796576_Nonlinear_time_series_methods_for_analyzing_behavioral_sequences}
                         {{\color{gray!60} "Nonlinear time series methods for analyzing behavioral sequences"}}}
            }

    \end{center}
\end{frame}

\begin{frame}
    \frametitle{}
    \begin{center}
        Dimentions

        DB parameters (e.g. shared\_buffers) +
        Workload parameters (e.g. rate of incoming queries) +
        Performance results (e.g. query latency)
    \end{center}
\end{frame}

\begin{frame}
    \frametitle{}
    \begin{center}
        \begin{tikzpicture}
        \begin{axis}[
            xlabel = shared\_buffers,
            ylabel = queries rate,
            zlabel = query latency,
            xtick = {0},
            ytick = {0},
            ztick = {0},
            ticklabel style = {font = \tiny},
            xmajorgrids=false,
            ymajorgrids=false,
            zmajorgrids=false,
        ]
        %\addplot3[
            %surf,
        %]
        \addplot3 [surf, domain=0:360, samples=60]
            { sin(x)*sin(y) };]
        %coordinates {
        %(0,0,0) (0,1,0) (0,2,0)

        %(1,0,0) (1,1,0.6) (1,2,0.7)

        %(2,0,0) (2,1,0.7) (2,2,1.8)
        %};
        %{exp(-x^2-y^2)*x};
        \end{axis}
        \end{tikzpicture}
    \end{center}
\end{frame}

\begin{frame}
    \frametitle{}
    \begin{center}
        \copyrightbox[b]
            {\includegraphics[width=0.8\textwidth,center]{probability-distribution.png}}
            {Probability distribution}

    \end{center}
\end{frame}

\begin{frame}
    \frametitle{}
    \begin{center}
        Benchmarking = exploring

        known vs unknown

        common vs particular
    \end{center}
\end{frame}

\section{PostgreSQL specifics}

\begin{frame}
    \frametitle{}
    \begin{center}
        Misconfiguration, too small or too big

        shared\_buffers, work\_mem, max\_wal\_size, wal\_writer\_flush\_after, checkpoint\_timeout, checkpoint\_completion\_target, checkpoint\_flush\_after
    \end{center}
\end{frame}

\begin{frame}
    \frametitle{}
    \begin{center}
        Not configuring OS

        huge pages, dirty\_bytes, dirty\_background\_bytes, readahead, fs journal, priorities
    \end{center}
\end{frame}

\begin{frame}
    \frametitle{}
    \begin{center}
        Noise, VM, noisy neighbors. hwlat.
    \end{center}
\end{frame}

\begin{frame}[fragile]{}
    \frametitle{}
    \begin{center}
        How long does it take to see the effect?

        Take into account checkpoints, autovacuum, bloat

        \begin{minted}[fontsize=\large]{shell-session}
#autovacuum_naptime = 1min		# time between autovacuum runs
#autovacuum_vacuum_threshold = 50	# min number of row updates before
                    # vacuum
#autovacuum_vacuum_insert_threshold = 1000	# min number of row inserts
                    # before vacuum; -1 disables insert
                    # vacuums
#autovacuum_vacuum_scale_factor = 0.2	# fraction of table size before vacuum
#autovacuum_vacuum_insert_scale_factor = 0.2	# fraction of inserts over table
                    # size before insert vacuum
#autovacuum_vacuum_cost_delay = 2ms	# default vacuum cost delay for
                    # autovacuum, in milliseconds;
                    # -1 means use vacuum_cost_delay
#autovacuum_vacuum_cost_limit = -1	# default vacuum cost limit for
                    # autovacuum, -1 means use
                    # vacuum_cost_limit
        \end{minted}

    \end{center}
\end{frame}

\begin{frame}
    \frametitle{}
    \begin{center}
        Load generator performance. Open vs closed systems. Pgbench is closed,
        benchbase can be open. The example at the beginning is from pgbench
        logging vs not logging.
    \end{center}
\end{frame}

\begin{frame}
    \frametitle{}
    \begin{center}
        Hardware. fsync graph.
    \end{center}
\end{frame}

\section{Statistics}

\fontsize{18pt}{18}\selectfont
\begin{frame}
  \vspace*{2.5cm}
  \begin{minipage}[b][\paperheight]{\textwidth}
  \begin{center}

      %\raggedright%
      \linespread{1.0}%
      \usebeamerfont{title}%
      \usebeamercolor[fg]{title}%
      \if@noSmallCapitals%
        Questions?
      \else%
        \scshape{\color{black} Questions?}%
      \fi%
      \vspace*{0.3em}

      \usebeamerfont{subtitle}%
      \fontsize{13pt}{14}\selectfont
      \usebeamercolor[fg]{subtitle}%
        \begin{itemize}[label={}]
            \item {\color{black} \social\ @erthalion@fosstodon.org}
            \item {\color{black} \email\ ddolgov at redhat dot com}
        \end{itemize}
      \vspace*{2.5em}%

    \vfill
    \vspace*{2em}
  \end{center}
  \end{minipage}

\end{frame}

\end{document}
